% generated by roam.py handbook export -- не редактировать, править заметки
% 20260915183647-собственные_значения_положительно_определенных_симметричных_матриц.org (body 32f3a0ed8262)

\begin{theorem}[{Собственные значения положительно определенных, симметричных матриц}]\label[theorem]{note:025A8419-962C-4892-A4F8-FB496AAA938C}
  У \hyperref[note:771B93FA-654D-4233-9487-B4FDD356DDFE]{положительно определенной симметричной матрицы} все ее \hyperref[note:9D2905FC-8E68-4E6F-9243-364269F795B0]{собственные значения} являются положительными вещественными числами.
\end{theorem}

\begin{proof}
  Пусть $A$~--- положительно определенная симметричная матрица, $\lambda$~--- ее собственное значение. По \hyperref[note:4F4E5575-ABE6-4D9B-9C61-B76A3AA9701C]{теореме о симметричных матрицах} $\lambda \in \R$.

  Собственный вектор можно взять вещественным: матрица $A - \lambda I$~--- вещественная и вырожденная, значит у системы $(A - \lambda I) q = 0$ есть ненулевое решение $q \in \R^n$.

  Подставим $x = q$ в определение положительной определенности:
  \[
  0 < q^T A q = q^T (\lambda q) = \lambda q^T q = \lambda \norm{q}^2
  \]
  Поделим на $\norm{q}^2 > 0$ и получим $\lambda > 0$.
\end{proof}
